% ============================================================================
%  conclusion.tex — Lkhatima dyal TP
%  hna kankhlless dakchi lli t3ellemna
% ============================================================================

\chapter*{Conclusion}
\addcontentsline{toc}{chapter}{Conclusion}

\`{A} l'issue de ce troisi\`{e}me travail pratique, nous avons pu asseoir nos connaissances fondamentales en programmation orient\'{e}e objet avec Java. La manipulation des classes telles que \texttt{Point}, \texttt{CompteBancaire} ou \texttt{Produit} nous a permis de comprendre l'importance de l'encapsulation, l'utilit\'{e} des constructeurs pour l'initialisation des objets, ainsi que le d\'{e}coupage logique des fonctionnalit\'{e}s en m\'{e}thodes distinctes. En outre, la cr\'{e}ation d'exercices interactifs tels que le \texttt{JeuDevinette} a mis en \'{e}vidence l'interaction entre les entr\'{e}es utilisateurs (avec la classe Scanner) et la logique interne des objets.\\

Ces concepts constituent d\'{e}sormais une base solide pour aborder la suite de notre apprentissage en Java, notamment pour les notions avanc\'{e}es d'h\'{e}ritage, de polymorphisme et surtout la structuration de projets logiciels complexes r\'{e}utilisables.
